% Slides 4--9: equations, equilibria, stability, persistence, and bifurcation.

\begin{frame}{The Dimensionless Model}
  \begin{columns}[T,onlytextwidth]
    \begin{column}{0.62\textwidth}
      \begin{align*}
        \dot{\Sstate} &= 1-\Sstate-f_1(\Sstate)\Xstate,\\
        \dot{\Xstate} &= \Xstate\bigl(f_1(\Sstate)-D_1\bigr)-f_2(\Xstate)\Ystate,\\
        \dot{\Ystate} &= \Ystate\bigl(f_2(\Xstate)-D_2\bigr)-f_3(\Ystate)\Zstate,\\
        \dot{\Zstate} &= \Zstate\bigl(f_3(\Ystate)-D_3\bigr).
      \end{align*}
    \end{column}
    \begin{column}{0.34\textwidth}
      \keypoint{Normalization}
      \begin{itemize}
        \item Nutrient input: $1$
        \item Nutrient dilution: $-S$
        \item Time scaled by dilution
        \item State $X=(S,x,y,z)^\mathsf{T}$
      \end{itemize}
      \vspace{0.2cm}
      \keypoint{Interpretation}
      \begin{itemize}
        \item Positive terms create biomass.
        \item $D_i$ removes biomass.
        \item Predation transfers biomass upward.
      \end{itemize}
    \end{column}
  \end{columns}
\end{frame}

\begin{frame}{Monod Growth and Parameters}
  \begin{columns}[T,onlytextwidth]
    \begin{column}{0.54\textwidth}
      \vspace{0.15cm}
      \[
        f_i(u)=\frac{a_i u}{k_i+u},\qquad
        f_i'(u)=\frac{a_i k_i}{(k_i+u)^2}
      \]
      \[
        f_i^{-1}(d)=\frac{d k_i}{a_i-d},
        \qquad 0\le d<a_i.
      \]
      \vspace{0.25cm}
      \begin{block}{Why Monod functions?}
        They are increasing, vanish without food, saturate at $a_i$, and
        preserve the paper's assumptions while giving reproducible formulas.
      \end{block}
    \end{column}
    \begin{column}{0.42\textwidth}
      {\fontsize{11}{14}\selectfont
      \keypoint{Parameter meanings}\par
      $a_i$: maximum growth rate\\
      $k_i$: half-saturation level\\
      $D_1$: prey removal rate\\
      $D_2$: first-predator removal rate\\
      $D_3$: top-predator removal rate
      \vspace{0.35cm}

      \keypoint{Simulation input}
      \[
        (a_1,k_1,a_2,k_2,a_3,k_3,D_1,D_2,D_3)
      \]}
    \end{column}
  \end{columns}
\end{frame}

\begin{frame}{Feasible Dynamics}
  \begin{columns}[T,onlytextwidth]
    \begin{column}{0.51\textwidth}
      \keypoint{Nonnegative invariant region}
      \begin{itemize}
        \item At $S=0$, $\dot S=1>0$.
        \item At $x=0$, $y=0$, or $z=0$, the corresponding derivative is zero.
        \item A positive trajectory cannot cross a coordinate plane.
      \end{itemize}
      \vspace{0.25cm}
      \keypoint{Total concentration}
      \[
        T=S+x+y+z,
      \]
      \[
        \dot T=1-S-D_1x-D_2y-D_3z.
      \]
    \end{column}
    \begin{column}{0.45\textwidth}
      Let
      \[
        D_{\min}=\min(1,D_1,D_2,D_3),
      \]
      \[
        D_{\max}=\max(1,D_1,D_2,D_3).
      \]
      Then
      \[
        1-D_{\max}T\le \dot T\le 1-D_{\min}T.
      \]
      \keypoint{Consequence}\par
      Solutions stay meaningful and eventually enter a common bounded region.
      The model is dissipative.
    \end{column}
  \end{columns}
\end{frame}

\begin{frame}{Four Equilibria, Four Regimes}
  \renewcommand{\arraystretch}{1.18}
  \begin{tabular}{clll}
    \toprule
    State & Positive populations & Defining relations & Biological regime \\
    \midrule
    $P_0$ & none
      & $(1,0,0,0)$
      & Complete washout \\
    $P_1$ & $x$
      & $f_1(S_1)=D_1$
      & Prey only \\
    $P_2$ & $x,y$
      & $f_2(x_2)=D_2$
      & Top predator absent \\
    $P_3$ & $x,y,z$
      & $f_3(y_3)=D_3$
      & Full coexistence \\
    \bottomrule
  \end{tabular}

  \vspace{0.45cm}
  \begin{columns}[T,onlytextwidth]
    \begin{column}{0.48\textwidth}
      \[
        S_1=f_1^{-1}(D_1),\qquad
        x_1=\frac{1-S_1}{D_1}
      \]
    \end{column}
    \begin{column}{0.48\textwidth}
      \[
        x_2=f_2^{-1}(D_2),\qquad
        y_3=f_3^{-1}(D_3)
      \]
    \end{column}
  \end{columns}
  \vspace{0.15cm}
  \centering\color{MidGray}
  $P_3$ requires a coupled numerical solve for $S_3$ and $x_3$.
\end{frame}

\begin{frame}{Invasion Controls Stability}
  \begin{columns}[T,onlytextwidth]
    \begin{column}{0.55\textwidth}
      \keypoint{Per-capita invasion rates}
      \begin{align*}
        \lambda_x &= f_1(1)-D_1,\\
        \lambda_y &= f_2(x_1)-D_2,\\
        \lambda_z &= f_3(y_2)-D_3.
      \end{align*}
      Positive invasion creates the next trophic regime; negative invasion
      stabilizes the lower-level boundary equilibrium.
    \end{column}
    \begin{column}{0.41\textwidth}
      {\fontsize{11}{14}\selectfont
      \renewcommand{\arraystretch}{1.25}
      \begin{tabular}{cl}
        \toprule
        Equilibrium & Main stability test \\
        \midrule
        $P_0$ & $\lambda_x<0$ \\
        $P_1$ & $\lambda_y<0$ \\
        $P_2$ & cubic RH and $\lambda_z<0$ \\
        $P_3$ & quartic RH test \\
        \bottomrule
      \end{tabular}
      \vspace{0.35cm}
      \begin{block}{Local stability}
        Every eigenvalue of the Jacobian must have negative real part.
      \end{block}}
    \end{column}
  \end{columns}
\end{frame}

\begin{frame}{Persistence, Hopf, and Model Choice}
  \begin{columns}[T,onlytextwidth]
    \begin{column}{0.47\textwidth}
      \keypoint{Uniform persistence}
      \begin{itemize}
        \item Dissipativity prevents unbounded escape.
        \item Boundary equilibria must repel each missing predator.
        \item Under the paper's extra assumptions, every positive component
              stays bounded away from zero asymptotically.
      \end{itemize}
      \[
        \liminf_{t\to\infty}(S,x,y,z)\ge \eta>0.
      \]
    \end{column}
    \begin{column}{0.49\textwidth}
      \keypoint{Hopf boundaries}
      \[
        P_2:\quad A_1A_2=A_3
      \]
      \[
        P_3:\quad C_1C_2C_3=C_3^2+C_1^2C_4
      \]
      A conjugate eigenvalue pair must cross the imaginary axis with nonzero
      speed; damped oscillation alone is not a proof of bifurcation.
    \end{column}
  \end{columns}
  \vspace{0.1cm}
  \keypoint{Implementation choice:} printed equation (1.2) conflicts with the
  later Jacobian, so the code follows the internally consistent analytical form.
\end{frame}
